\documentclass[12pt, a4paper]{article}
\usepackage{uepb_relatorio_config}

% --- PACOTES E MACROS ADICIONAIS ---
\usepackage{amsmath, amssymb}
\usepackage{float}
\usepackage{paineis_tikz}

% Referência à Aula2 (Prof. Robson Pequeno de Sousa) com nota de rodapé
% \aula: primeira menção, com a nota completa; \aularef: menções seguintes, que remetem a essa nota.
\newcommand{\aula}{Aula~2\footnote{\label{fn:aula}SOUSA, Robson Pequeno de. \textit{Sistemas de coordenadas}: Aula2. Material da disciplina de Computação Gráfica. As referências bibliográficas constam ao final do material do professor. Ver também \cite{aula2}.}}
\newcommand{\aularef}{Aula~2\footnote{Ver a nota~\ref{fn:aula}.}}


\begin{document}

\emergencystretch=3em
\hbadness=10000
\vbadness=10000
\flushbottom
\setlength{\skip\footins}{12pt plus 4pt minus 4pt}
\setlength{\footnotesep}{11pt}

% --- CAPA ---
\logomarca{imagens/marca-uepb.png}
\universidade{UNIVERSIDADE ESTADUAL DA PARAÍBA - UEPB}
\centro{CENTRO DE CIÊNCIAS E TECNOLOGIA - CCT}
\departamento{DEPARTAMENTO DE COMPUTAÇÃO}
\curso{CURSO DE CIÊNCIA DA COMPUTAÇÃO}

\autores{
    Antonio Wenicios Ademar da Silva \\
    Julia Beatriz Pamplona de Sousa \\
    Julio Pedro da Silva \\
    Livia Sousa Silva Ferreira    \\
    Thales Henrique de Souza Teles \\
}

\titulo{TRANSFORMAÇÕES DE COORDENADAS}

\cidadeuf{Campina Grande - PB}
\dataano{2026}

\imprimircapa

\tableofcontents
\newpage

\section{Introdução}
\label{sec:introducao}
O presente relatório descreve o desenvolvimento do primeiro projeto prático da disciplina de Computação Gráfica. O trabalho aborda um dos conceitos fundamentais da computação gráfica bidimensional: a transformação entre diferentes sistemas de coordenadas utilizados durante o processo de visualização.

Nesse processo, um ponto inicialmente definido no Sistema de Coordenadas do Mundo é transformado para o Sistema de Coordenadas Normalizadas do Dispositivo (NDC) e, posteriormente, para as Coordenadas do Dispositivo, que correspondem à posição do ponto na tela em pixels.

O projeto consiste no desenvolvimento de uma aplicação em Java capaz de realizar e demonstrar essas transformações matemáticas. A aplicação utiliza a biblioteca Java Swing para apresentar os resultados graficamente, permitindo visualizar o mapeamento de um ponto do espaço do mundo até sua representação na tela.

\section{Objetivos}
\label{sec:objetivos}

\subsection{Objetivo geral}
Compreender e implementar, na prática, as operações matemáticas necessárias para transformar coordenadas bidimensionais entre o Sistema de Coordenadas do Mundo, o Sistema de Coordenadas Normalizadas do Dispositivo (NDC) e as Coordenadas do Dispositivo, utilizando como representação visual um pixel de cor verde na tela.

\subsection{Objetivos específicos}
\begin{itemize}
    \item \textbf{Modelagem de dados:} Estruturar as classes necessárias para representar pontos cartesianos e os limites das janelas de visualização.
    \item \textbf{Normalização:} Desenvolver os algoritmos de conversão entre o Sistema de Coordenadas do Mundo e o Sistema de Coordenadas Normalizadas do Dispositivo (NDC), considerando os intervalos $[0,1]$ e $[-1,1]$.
    \item \textbf{Mapeamento para a tela:} Implementar a conversão das coordenadas NDC para as Coordenadas do Dispositivo, transformando as posições normalizadas em coordenadas de pixels de acordo com as dimensões da tela.
    \item \textbf{Aplicação interativa:} Construir uma interface gráfica utilizando a biblioteca Java Swing, permitindo ao usuário visualizar e validar o funcionamento das transformações implementadas.
\end{itemize}

\section{Fundamentação teórica}
\label{sec:fundamentacao-teorica}
Para o desenvolvimento do rasterizador e o desenho do pixel na tela, diferentes espaços matemáticos precisam ser compreendidos. Nesta seção, cada sistema de coordenadas é apresentado junto de um gráfico que posiciona um mesmo ponto de exemplo e, em seguida, as fórmulas de transformação são deduzidas passo a passo. A referência principal para as expressões utilizadas é a \aula, e as deduções aqui apresentadas buscam chegar às mesmas fórmulas dessa aula; a correspondência de notação é apresentada na Seção~\ref{subsec:sintese}.

\subsection{Sistemas de coordenadas}
Os sistemas de coordenadas definem espaços onde pontos e objetos são posicionados e manipulados. Na Computação Gráfica, utilizamos sistemas abstratos e sistemas físicos para garantir que a imagem processada seja desenhada de forma correta e independente da resolução de tela.

Neste trabalho são utilizados três sistemas: o sistema de coordenadas do \textbf{mundo} (definido pelo usuário), o sistema de coordenadas \textbf{normalizadas do dispositivo} (NDC), e o sistema de \textbf{coordenadas de dispositivo} (DC). Um ponto especificado no mundo percorre o caminho
\[
    \text{Mundo} \longrightarrow \text{NDC} \longrightarrow \text{DC},
\]
e o caminho inverso é percorrido quando uma posição da tela (por exemplo, um clique do mouse) é convertida de volta para o mundo. A etapa intermediária (NDC) é estudada em dois cenários: $[0,1] \times [0,1]$ e $[-1,1] \times [-1,1]$.

Um ponto é denotado por $P_w = (x_w, y_w)$ no mundo, $P_n = (x_n, y_n)$ no NDC e $P_d = (x_d, y_d)$ no dispositivo. A razão entre o deslocamento de um ponto e a largura total do sistema é denotada pela letra $R$; o primeiro índice de $R$ indica o sistema (\emph{w}, \emph{n} ou \emph{d}) e o segundo, o eixo.

Todos os gráficos deste relatório seguem a convenção adotada em sala na \aularef, em que o mundo, o NDC e o dispositivo foram desenhados no quadro com a origem (o mínimo de cada eixo) no canto inferior esquerdo e o eixo $y$ crescendo para cima \cite{quadro2}, convenção também usada em \cite{foley}. As transformações deduzidas são apenas escalas por razões e não invertem nenhum eixo; por isso, o mesmo ponto ocupa a mesma posição relativa em todos os sistemas. A orientação do display da implementação em Java, que difere dessa convenção, é discutida na Seção~\ref{sec:metodologia}.

\subsection{Coordenadas de dispositivo}
O sistema de Coordenadas de Dispositivo (DC) é um sistema \textbf{discreto} e representa a matriz física de pixels de um monitor ou tela, possuindo resolução finita com $W$ pixels na horizontal e $H$ pixels na vertical. Os pixels são endereçados por dois números inteiros, as coordenadas horizontal e vertical. Adota-se a origem $(0,0)$ no canto inferior esquerdo, com $y_d$ crescendo para cima (Figura~\ref{fig:dc}), de modo que qualquer ponto $P_d = (x_d, y_d)$ é um par de números inteiros tal que
\begin{equation}
    x_d \in \{0, 1, \dots, W-1\}, \qquad y_d \in \{0, 1, \dots, H-1\}.
\end{equation}
Como a contagem dos pixels começa em $0$, o último pixel de cada eixo possui índice $W-1$ (horizontal) ou $H-1$ (vertical). Assim, a largura total da tela, medida em índices de pixel, é $W-1$ e a altura total é $H-1$.

Para localizar $P_d$ dentro da tela, medem-se os deslocamentos a partir da origem $(0,0)$ (Figura~\ref{fig:dc}):
\begin{equation}
    \Delta x_d = x_d - 0, \qquad \Delta y_d = y_d - 0.
\end{equation}
A razão entre o deslocamento $\Delta x_d$ e a largura total da tela é
\begin{equation}
    R_{dx} = \frac{\Delta x_d}{W-1} = \frac{x_d}{W-1},
    \label{eq:razao-dc-x}
\end{equation}
isto é, a fração da largura da tela já percorrida até o pixel. Analogamente para $y$:
\begin{equation}
    R_{dy} = \frac{\Delta y_d}{H-1} = \frac{y_d}{H-1}.
    \label{eq:razao-dc-y}
\end{equation}
Tem-se $R_{dx}, R_{dy} \in [0,1]$.

\begin{figure}[H]
    \centering
    \begin{tikzpicture}
        \resetpainel \presetdc \def\dims{1}\def\ptit{}
        \painel{dc}
    \end{tikzpicture}
    \caption{Sistema de coordenadas de dispositivo (discreto) com um ponto de exemplo $P_d = (x_d, y_d)$. As cotas em azul indicam o deslocamento do ponto e a largura total em cada eixo.}
    \label{fig:dc}
\end{figure}

\subsection{Janela de mundo}
A Janela de Mundo é a região do universo contínuo onde os objetos virtuais estão definidos. Suas fronteiras são parametrizadas pelas coordenadas reais $x_{wmin}$, $x_{wmax}$, $y_{wmin}$ e $y_{wmax}$. Qualquer coordenada do mundo é definida pelo par contínuo $P_w = (x_w, y_w)$.

Para localizar $P_w$ dentro da janela, medem-se os deslocamentos a partir do canto inferior esquerdo $(x_{wmin}, y_{wmin})$ (Figura~\ref{fig:mundo}):
\begin{equation}
    \Delta x_w = x_w - x_{wmin}, \qquad \Delta y_w = y_w - y_{wmin}.
\end{equation}
A razão entre o deslocamento $\Delta x_w$ e a largura total da janela é
\begin{equation}
    R_{wx} = \frac{\Delta x_w}{x_{wmax} - x_{wmin}} = \frac{x_w - x_{wmin}}{x_{wmax} - x_{wmin}},
    \label{eq:razao-mundo-x}
\end{equation}
que responde à pergunta: ``quantas vezes $\Delta x_w$ cabe na largura total?''. Analogamente para $y$:
\begin{equation}
    R_{wy} = \frac{\Delta y_w}{y_{wmax} - y_{wmin}} = \frac{y_w - y_{wmin}}{y_{wmax} - y_{wmin}}.
    \label{eq:razao-mundo-y}
\end{equation}
Para um ponto dentro da janela, tem-se $R_{wx}, R_{wy} \in [0,1]$; valores fora desse intervalo indicam um ponto fora da janela de mundo.

\begin{figure}[H]
    \centering
    \begin{tikzpicture}
        \resetpainel \presetmundo \def\dims{1}\def\ptit{}
        \painel{mundo}
    \end{tikzpicture}
    \caption{Sistema de coordenadas do mundo com um ponto de exemplo $P_w = (x_w, y_w)$, os deslocamentos $\Delta x_w$ e $\Delta y_w$ e as larguras totais da janela.}
    \label{fig:mundo}
\end{figure}

\subsection{NDC [0,1] e NDC [-1,1]}
Para contornar o problema de resoluções variadas entre dispositivos, adota-se um sistema intermediário independente de resolução, o NDC (\textit{Normalized Device Coordinates}).
\begin{itemize}
    \item \textbf{Cenário (i) - NDC [0,1]:} As coordenadas são normalizadas para flutuarem entre $0$ (mínimo) e $1$ (máximo), com a origem no canto inferior esquerdo.
    \item \textbf{Cenário (ii) - NDC [-1,1]:} Normalização comum no pipeline padrão do OpenGL, onde a origem $(0,0)$ reside no centro lógico da tela, estendendo-se para os extremos de $-1$ a $1$.
\end{itemize}
Os pontos NDC são números reais, denotados por $P_n = (x_n, y_n)$; quando for necessário distinguir os cenários, escreve-se $P_n = (x_n^{(i)}, y_n^{(i)})$ no cenário (i) e $P_n = (x_n^{(ii)}, y_n^{(ii)})$ no cenário (ii).

Assim como nos demais sistemas, $P_n$ é localizado pelos deslocamentos medidos a partir do extremo inferior/esquerdo de cada eixo (Figura~\ref{fig:ndc}):
\begin{equation}
    \Delta x_n = x_n - x_{nmin}, \qquad \Delta y_n = y_n - y_{nmin},
\end{equation}
sendo $x_{nmin} = y_{nmin} = 0$ no cenário (i) e $x_{nmin} = y_{nmin} = -1$ no cenário (ii); em ambos, $x_{nmax} = y_{nmax} = 1$.

No cenário (i), o deslocamento vale $\Delta x_n^{(i)} = x_n^{(i)} - 0$ e a largura total vale $1 - 0 = 1$, logo
\begin{equation}
    R_{nx}^{(i)} = \frac{x_n^{(i)} - 0}{1 - 0} = x_n^{(i)}.
    \label{eq:razao-ndc01-x}
\end{equation}
No cenário (ii), o deslocamento vale $\Delta x_n^{(ii)} = x_n^{(ii)} - (-1) = x_n^{(ii)} + 1$ e a largura total vale $1 - (-1) = 2$, logo
\begin{equation}
    R_{nx}^{(ii)} = \frac{x_n^{(ii)} - (-1)}{1 - (-1)} = \frac{x_n^{(ii)} + 1}{2}.
    \label{eq:razao-ndc11-x}
\end{equation}
Os resultados para o eixo $y$ são análogos: $R_{ny}^{(i)} = y_n^{(i)}$ e $R_{ny}^{(ii)} = \frac{y_n^{(ii)} + 1}{2}$.

\begin{figure}[H]
    \centering
    \begin{tikzpicture}
        \begin{scope}
            \resetpainel \presetndci \def\dims{1}\def\pw{4.0}
            \painel{ndcA}
        \end{scope}

        \begin{scope}[xshift=7.4cm]
            \resetpainel \presetndcii \def\dims{1}\def\pw{4.0}
            \painel{ndcB}
        \end{scope}
    \end{tikzpicture}
    \caption{Sistemas de coordenadas normalizadas do dispositivo nos cenários (i) e (ii), com um ponto de exemplo $P_n$. Em ambos, os deslocamentos são medidos a partir do extremo inferior/esquerdo de cada eixo.}
    \label{fig:ndc}
\end{figure}

\subsection{Transformações entre os sistemas}
Todas as conversões seguem a mesma ideia. Em qualquer sistema, a posição de um ponto ao longo de um eixo é descrita pela razão
\[
    R = \frac{\text{``deslocamento a partir do mínimo''}}{\text{``largura total''}} = \frac{q - q_{min}}{q_{max} - q_{min}}.
\]
em que $q$ denota, de forma genérica, a coordenada de um ponto ao longo de um dos eixos ($x$ ou $y$) em qualquer um dos três sistemas (Mundo, NDC ou DC), e $q_{min}$, $q_{max}$ são os extremos desse mesmo eixo nesse sistema. Por exemplo, $q = x_w$, $q_{min} = x_{wmin}$ e $q_{max} = x_{wmax}$ no caso do eixo $x$ no Mundo. Como as transformações apenas escalam e transladam cada eixo, o ponto ocupa a \textbf{mesma proporção} da largura total em todos os sistemas. Basta, portanto, \textbf{igualar as razões} dos dois sistemas e isolar a incógnita, procedimento seguido em sala no quadro \cite{quadro2}. Considera-se $x_{wmin} < x_{wmax}$ e $y_{wmin} < y_{wmax}$. Cada transformação é deduzida para os dois cenários de NDC e o nome do procedimento correspondente aparece entre parênteses no título. As deduções são feitas para $x$; as expressões finais incluem também $y$, obtido de forma análoga.

% ------------------------------------------------------------
\subsubsection{\texorpdfstring{Mundo $\rightarrow$ NDC}{Mundo -> NDC} (\texttt{user\_to\_ndc})}
Dado $P_w = (x_w, y_w)$ no mundo, obtém-se o ponto $P_n = (x_n, y_n)$ no NDC.

\paragraph{Cenário (i) --- NDC [0,1].} Igualando a razão no mundo, $R_{wx}$, à razão no NDC [0,1], $R_{nx}^{(i)}$ (Figura~\ref{fig:mundo-ndc01}):
\[
    \frac{x_w - x_{wmin}}{x_{wmax} - x_{wmin}} = \frac{x_n^{(i)} - 0}{1 - 0} = x_n^{(i)}.
\]
O termo $x_n^{(i)}$ já está isolado. Fazendo o mesmo para $y$:
\begin{equation}
    \boxed{\begin{aligned}
        x_n^{(i)} &= \frac{x_w - x_{wmin}}{x_{wmax} - x_{wmin}} \\[2mm]
        y_n^{(i)} &= \frac{y_w - y_{wmin}}{y_{wmax} - y_{wmin}}
    \end{aligned}}
    \label{eq:user-ndc-i}
\end{equation}

\begin{figure}[H]
    \centering
    \mapafig{\presetmundo}{\presetndci}
    \caption{Conversão nos dois sentidos entre o sistema do mundo e o NDC $[0,1] \times [0,1]$: \texttt{user\_to\_ndc} ($\rightarrow$) e \texttt{ndc\_to\_user} ($\leftarrow$).}
    \label{fig:mundo-ndc01}
\end{figure}

\paragraph{Cenário (ii) --- NDC [-1,1].} O extremo inferior do eixo é $-1$, então o deslocamento é $x_n^{(ii)} - (-1)$ e a largura total é $1 - (-1) = 2$ (Figura~\ref{fig:mundo-ndc11}). Igualando as razões:
\[
    \frac{x_w - x_{wmin}}{x_{wmax} - x_{wmin}} = \frac{x_n^{(ii)} + 1}{2}.
\]
Multiplicando ambos os lados por $2$ e subtraindo $1$:
\[
    x_n^{(ii)} = 2 \cdot \frac{x_w - x_{wmin}}{x_{wmax} - x_{wmin}} - 1
    = \frac{2x_w - x_{wmax} - x_{wmin}}{x_{wmax} - x_{wmin}}.
\]
Fazendo o mesmo para $y$:
\begin{equation}
    \boxed{\begin{aligned}
        x_n^{(ii)} &= 2 \cdot \frac{x_w - x_{wmin}}{x_{wmax} - x_{wmin}} - 1 \\[2mm]
        y_n^{(ii)} &= 2 \cdot \frac{y_w - y_{wmin}}{y_{wmax} - y_{wmin}} - 1
    \end{aligned}}
    \label{eq:user-ndc-ii}
\end{equation}
Note que $x_n^{(ii)} = 2\,x_n^{(i)} - 1$: as bordas $x_{wmin}$ e $x_{wmax}$ e o centro da janela são levados a $-1$, $1$ e $0$, respectivamente.

\begin{figure}[H]
    \centering
    \mapafig{\presetmundo}{\presetndcii}
    \caption{Conversão nos dois sentidos entre o sistema do mundo e o NDC $[-1,1] \times [-1,1]$: \texttt{user\_to\_ndc} ($\rightarrow$) e \texttt{ndc\_to\_user} ($\leftarrow$).}
    \label{fig:mundo-ndc11}
\end{figure}

% ------------------------------------------------------------
\subsubsection{\texorpdfstring{NDC $\rightarrow$ Mundo}{NDC -> Mundo} (\texttt{ndc\_to\_user})}
Dado $P_n = (x_n, y_n)$ no NDC, obtém-se $P_w = (x_w, y_w)$ no mundo (transformação inversa, Figuras~\ref{fig:mundo-ndc01} e~\ref{fig:mundo-ndc11}). Parte-se da mesma igualdade de razões, agora isolando $x_w$.

\paragraph{Cenário (i) --- NDC [0,1].} Da igualdade $\dfrac{x_w - x_{wmin}}{x_{wmax} - x_{wmin}} = x_n^{(i)}$, multiplica-se por $x_{wmax} - x_{wmin}$ e soma-se $x_{wmin}$:
\begin{equation}
    \boxed{\begin{aligned}
        x_w &= x_{wmin} + x_n^{(i)} \, (x_{wmax} - x_{wmin}) \\
        y_w &= y_{wmin} + y_n^{(i)} \, (y_{wmax} - y_{wmin})
    \end{aligned}}
    \label{eq:ndc-user-i}
\end{equation}

\paragraph{Cenário (ii) --- NDC [-1,1].} Da igualdade $\dfrac{x_w - x_{wmin}}{x_{wmax} - x_{wmin}} = \dfrac{x_n^{(ii)} + 1}{2}$, procede-se da mesma forma:
\begin{equation}
    \boxed{\begin{aligned}
        x_w &= x_{wmin} + \frac{x_n^{(ii)} + 1}{2} \, (x_{wmax} - x_{wmin}) \\[2mm]
        y_w &= y_{wmin} + \frac{y_n^{(ii)} + 1}{2} \, (y_{wmax} - y_{wmin})
    \end{aligned}}
    \label{eq:ndc-user-ii}
\end{equation}

% ------------------------------------------------------------
\subsubsection{\texorpdfstring{NDC $\rightarrow$ DC}{NDC -> DC} (\texttt{ndc\_to\_dc})}
Dado $P_n = (x_n, y_n)$ no NDC, obtém-se o ponto equivalente $P_d = (x_d, y_d)$ nas coordenadas de dispositivo. A razão no dispositivo é $R_{dx} = \dfrac{x_d}{W-1}$.

\paragraph{Cenário (i) --- NDC [0,1].} Igualando $R_{nx}^{(i)} = R_{dx}$ (Figura~\ref{fig:ndc01-dc}):
\[
    \frac{x_n^{(i)} - 0}{1 - 0} = \frac{x_d}{W-1}
    \quad \Rightarrow \quad
    x_d = x_n^{(i)} \, (W-1).
\]
Como $x_n^{(i)} \in [0,1] \subset \mathbb{R}$, esse produto em geral não é inteiro, mas o índice de um pixel precisa ser. Realiza-se, então, um arredondamento para o inteiro mais próximo, definido por
\begin{equation}
    \operatorname{round}(x) =
    \begin{cases}
        \lceil x \rceil, & \text{se } x - \lfloor x \rfloor \ge 0{,}5 \\
        \lfloor x \rfloor, & \text{se } x - \lfloor x \rfloor < 0{,}5
    \end{cases}
    \label{eq:round}
\end{equation}
Analogamente para $y$, com $H-1$:
\begin{equation}
    \boxed{\begin{aligned}
        x_d &= \operatorname{round}\!\left( x_n^{(i)} \, (W-1) \right) \\
        y_d &= \operatorname{round}\!\left( y_n^{(i)} \, (H-1) \right)
    \end{aligned}}
    \label{eq:ndc-dc-i}
\end{equation}

O trecho a seguir, em C, ilustra apenas o princípio do arredondamento: soma-se $0{,}5$ ao valor e trunca-se o resultado para inteiro\footnote{Em C, o \textit{cast} \texttt{(int)} converte um valor de ponto flutuante para inteiro descartando a parte decimal (truncamento em direção a zero), mantendo apenas a parte inteira do número. Por exemplo, \texttt{(int)(599.25 + 0.5)} = \texttt{(int)(599.75)} = \texttt{599}.}; como os valores tratados aqui são sempre não negativos, essa truncagem equivale a tomar o piso, reproduzindo exatamente a Definição~\eqref{eq:round}:
\begin{verbatim}
int round(float a) {
    return (int)(a + 0.5);
}
\end{verbatim}

No projeto\footnote{Código-fonte disponível no repositório do projeto \cite{repo}.}, a função efetivamente utilizada foi a \texttt{Math.round()} da biblioteca padrão do Java, que segue o mesmo princípio: soma-se $0{,}5$ ao valor de ponto flutuante e toma-se o piso do resultado para obter o inteiro mais próximo.
O arredondamento limita o erro de posição a meio pixel, enquanto o piso ($\lfloor \cdot \rfloor$) poderia deslocar o ponto em quase um pixel inteiro.

\begin{figure}[H]
    \centering
    \mapafig{\presetndci}{\presetdc}
    \caption{Conversão nos dois sentidos entre o NDC $[0,1] \times [0,1]$ e o dispositivo: \texttt{ndc\_to\_dc} ($\rightarrow$) e \texttt{inp\_to\_ndc} ($\leftarrow$).}
    \label{fig:ndc01-dc}
\end{figure}

\paragraph{Cenário (ii) --- NDC [-1,1].} Igualando $R_{nx}^{(ii)} = R_{dx}$ (Figura~\ref{fig:ndc11-dc}) e multiplicando por $W-1$:
\[
    \frac{x_n^{(ii)} + 1}{2} = \frac{x_d}{W-1}
    \quad \Rightarrow \quad
    x_d = \frac{x_n^{(ii)} + 1}{2} \, (W-1).
\]
O fator $\left(x_n^{(ii)} + 1\right)/2$ pertence a $[0,1]$ e o resultado, em geral, não é inteiro; aplica-se a função $\operatorname{round}$ da Definição~\eqref{eq:round}:
\begin{equation}
    \boxed{\begin{aligned}
        x_d &= \operatorname{round}\!\left( \frac{x_n^{(ii)} + 1}{2} \, (W-1) \right) \\[2mm]
        y_d &= \operatorname{round}\!\left( \frac{y_n^{(ii)} + 1}{2} \, (H-1) \right)
    \end{aligned}}
    \label{eq:ndc-dc-ii}
\end{equation}

\begin{figure}[H]
    \centering
    \mapafig{\presetndcii}{\presetdc}
    \caption{Conversão nos dois sentidos entre o NDC $[-1,1] \times [-1,1]$ e o dispositivo: \texttt{ndc\_to\_dc} ($\rightarrow$) e \texttt{inp\_to\_ndc} ($\leftarrow$).}
    \label{fig:ndc11-dc}
\end{figure}

% ------------------------------------------------------------
\subsubsection{\texorpdfstring{DC $\rightarrow$ NDC}{DC -> NDC} (\texttt{inp\_to\_ndc})}
Dado $P_d = (x_d, y_d)$ em pixels, por exemplo a posição de um clique do mouse, obtém-se $P_n = (x_n, y_n)$ no NDC (transformação inversa, Figuras~\ref{fig:ndc01-dc} e~\ref{fig:ndc11-dc}). Parte-se da mesma igualdade $R_{dx} = R_{nx}$, agora isolando $x_n$. Como $x_d$ é inteiro e $x_n$ é real, não há arredondamento, mas a divisão deve ser feita em ponto flutuante.

\paragraph{Cenário (i) --- NDC [0,1].} Da igualdade $\dfrac{x_d}{W-1} = \dfrac{x_n^{(i)} - 0}{1 - 0}$:
\begin{equation}
    \boxed{\begin{aligned}
        x_n^{(i)} &= \frac{x_d}{W-1} \\[2mm]
        y_n^{(i)} &= \frac{y_d}{H-1}
    \end{aligned}}
    \label{eq:dc-ndc-i}
\end{equation}

\paragraph{Cenário (ii) --- NDC [-1,1].} Da igualdade $\dfrac{x_d}{W-1} = \dfrac{x_n^{(ii)} + 1}{2}$, multiplica-se por $2$ e subtrai-se $1$:
\begin{equation}
    \boxed{\begin{aligned}
        x_n^{(ii)} &= 2 \cdot \frac{x_d}{W-1} - 1 \\[2mm]
        y_n^{(ii)} &= 2 \cdot \frac{y_d}{H-1} - 1
    \end{aligned}}
    \label{eq:dc-ndc-ii}
\end{equation}

% ------------------------------------------------------------
\subsubsection{Síntese e composição dos procedimentos}
\label{subsec:sintese}
A Tabela~\ref{tab:sintese} resume as expressões para a coordenada $x$; as de $y$ são obtidas trocando $x$ por $y$, $x_{wmin}$ e $x_{wmax}$ por $y_{wmin}$ e $y_{wmax}$, e $W$ por $H$.

\begin{table}[H]
    \centering
    \renewcommand{\arraystretch}{1.9}
    \small
    \begin{tabular}{l|c|l}
        \textbf{Procedimento} & \textbf{Cenário} & \textbf{Expressão para $x$} \\ \hline
        \texttt{user\_to\_ndc} & (i)  & $x_n^{(i)} = \dfrac{x_w - x_{wmin}}{x_{wmax} - x_{wmin}}$ \\
        \texttt{user\_to\_ndc} & (ii) & $x_n^{(ii)} = 2 \cdot \dfrac{x_w - x_{wmin}}{x_{wmax} - x_{wmin}} - 1$ \\
        \texttt{ndc\_to\_user} & (i)  & $x_w = x_{wmin} + x_n^{(i)} \, (x_{wmax} - x_{wmin})$ \\
        \texttt{ndc\_to\_user} & (ii) & $x_w = x_{wmin} + \dfrac{x_n^{(ii)} + 1}{2} \, (x_{wmax} - x_{wmin})$ \\
        \texttt{ndc\_to\_dc}   & (i)  & $x_d = \operatorname{round}\!\left( x_n^{(i)} \, (W-1) \right)$ \\
        \texttt{ndc\_to\_dc}   & (ii) & $x_d = \operatorname{round}\!\left( \dfrac{x_n^{(ii)} + 1}{2} \, (W-1) \right)$ \\
        \texttt{inp\_to\_ndc}  & (i)  & $x_n^{(i)} = \dfrac{x_d}{W-1}$ \\
        \texttt{inp\_to\_ndc}  & (ii) & $x_n^{(ii)} = 2 \cdot \dfrac{x_d}{W-1} - 1$
    \end{tabular}
    \caption{Resumo dos procedimentos de transformação de coordenadas.}
    \label{tab:sintese}
\end{table}

Compondo os procedimentos, e lembrando que $x_n^{(i)} = \left(x_n^{(ii)} + 1\right)/2 = R_{wx}$, os dois cenários levam às mesmas expressões diretas:
\begin{align*}
    \text{Mundo} \rightarrow \text{DC:} \quad & x_d = \operatorname{round}\!\left( \frac{x_w - x_{wmin}}{x_{wmax} - x_{wmin}} \, (W-1) \right), \\[1mm]
    \text{DC} \rightarrow \text{Mundo:} \quad & x_w = x_{wmin} + \frac{x_d}{W-1} \, (x_{wmax} - x_{wmin}).
\end{align*}
Isso mostra que a escolha entre NDC $[0,1]$ e NDC $[-1,1]$ altera apenas os valores intermediários, e não o pixel final nem o ponto recuperado no mundo, o que serve de critério de validação da implementação. Como o passo Mundo $\rightarrow$ DC arredonda para o pixel, a volta DC $\rightarrow$ Mundo recupera o ponto original apenas com erro de até meio pixel.

As expressões deduzidas coincidem com as apresentadas na \aularef, que adota o NDC $[0,1] \times [0,1]$ e escreve, para a transformação do mundo para o NDC e do NDC para o dispositivo,
\begin{align*}
    \mathit{ndcx} &= \frac{x - x_{min}}{x_{max} - x_{min}}, &
    \mathit{ndcy} &= \frac{y - y_{min}}{y_{max} - y_{min}}, \\[1mm]
    \mathit{dcx} &= \operatorname{round}\!\big(\mathit{ndcx}\,(\mathit{ndh} - 1)\big), &
    \mathit{dcy} &= \operatorname{round}\!\big(\mathit{ndcy}\,(\mathit{ndv} - 1)\big).
\end{align*}
A primeira linha corresponde à Equação~\eqref{eq:user-ndc-i} e a segunda à Equação~\eqref{eq:ndc-dc-i}, apenas com outra notação (Tabela~\ref{tab:notacao}). O cenário $[-1,1] \times [-1,1]$, solicitado no exercício da aula, é obtido pelas Equações~\eqref{eq:user-ndc-ii} e~\eqref{eq:ndc-dc-ii} e, como mostrado acima, leva à mesma expressão final para o pixel. Em ambas as fórmulas o eixo $y$ não sofre inversão: a aula não prevê reflexão entre o NDC e o dispositivo.

\begin{table}[H]
    \centering
    \renewcommand{\arraystretch}{1.5}
    \small
    \begin{tabular}{l|l}
        \textbf{Aula2} & \textbf{Este relatório} \\ \hline
        $x,\ y$ & $x_w,\ y_w$ \\
        $x_{min},\ x_{max},\ y_{min},\ y_{max}$ & $x_{wmin},\ x_{wmax},\ y_{wmin},\ y_{wmax}$ \\
        $\mathit{ndcx},\ \mathit{ndcy}$ & $x_n^{(i)},\ y_n^{(i)}$ \\
        $\mathit{dcx},\ \mathit{dcy}$ & $x_d,\ y_d$ \\
        $\mathit{ndh},\ \mathit{ndv}$ & $W,\ H$
    \end{tabular}
    \caption{Correspondência entre a notação da Aula2 e a deste relatório.}
    \label{tab:notacao}
\end{table}

% ------------------------------------------------------------
\subsubsection{Exemplo numérico}
Considere a janela de mundo $x_{wmin} = -2$, $x_{wmax} = 6$, $y_{wmin} = -1$, $y_{wmax} = 4$, um dispositivo de $W = 800$ por $H = 600$ pixels e o ponto $P_w = (4, 2)$.

\paragraph{Ida (Mundo $\rightarrow$ NDC $\rightarrow$ DC), cenário (i).} Pelas Equações~\eqref{eq:user-ndc-i} e~\eqref{eq:ndc-dc-i}:
\begin{align*}
    x_n^{(i)} &= \frac{4 - (-2)}{6 - (-2)} = 0{,}75, &
    y_n^{(i)} &= \frac{2 - (-1)}{4 - (-1)} = 0{,}6,
\end{align*}
\begin{align*}
    x_d &= \operatorname{round}(0{,}75 \cdot 799) = \operatorname{round}(599{,}25) = 599, \\
    y_d &= \operatorname{round}(0{,}6 \cdot 599) = \operatorname{round}(359{,}4) = 359.
\end{align*}

\begin{figure}[H]
    \centering
    \begin{tikzpicture}
        \def\ex{5.0}
        \begin{scope}
            \resetpainel \presetmundo
            \def\pw{3.3}\def\ph{2.3}\def\lfont{\scriptsize}
            \def\lxa{$-2$}\def\lxb{$6$}\def\lya{$-1$}\def\lyb{$4$}\def\lpt{$P_w = (4,\ 2)$}
            \painel{A}
        \end{scope}
        \begin{scope}[xshift=\ex cm]
            \resetpainel \presetndci
            \def\pw{3.3}\def\ph{2.3}\def\ptit{NDC $[0,1]^2$}
            \def\lpt{$P_n = (0{,}75;\ 0{,}6)$}
            \painel{B}
        \end{scope}
        \begin{scope}[xshift=2*\ex cm]
            \resetpainel \presetdc
            \def\pw{3.3}\def\ph{2.3}\def\ptit{DC}\def\lfont{\scriptsize}
            \def\lxb{$799$}\def\lyb{$599$}\def\lpt{$P_d = (599,\ 359)$}
            \painel{C}
        \end{scope}
        \draw[-{Stealth[length=2.5mm]}, thick, red!75!black] (A-P) -- (B-P);
        \draw[-{Stealth[length=2.5mm]}, thick, red!75!black] (B-P) -- (C-P);
    \end{tikzpicture}
    \caption{Exemplo numérico do cenário (i): o mesmo ponto no mundo, no NDC $[0,1]$ e no dispositivo. A ida percorre os sistemas da esquerda para a direita e a volta, da direita para a esquerda.}
    \label{fig:exemplo-i}
\end{figure}

\paragraph{Ida, cenário (ii).} Pelas Equações~\eqref{eq:user-ndc-ii} e~\eqref{eq:ndc-dc-ii}:
\begin{align*}
    x_n^{(ii)} &= 2 \cdot 0{,}75 - 1 = 0{,}5, &
    y_n^{(ii)} &= 2 \cdot 0{,}6 - 1 = 0{,}2, \\
    x_d &= \operatorname{round}\!\left( \tfrac{0{,}5 + 1}{2} \cdot 799 \right) = 599, &
    y_d &= \operatorname{round}\!\left( \tfrac{0{,}2 + 1}{2} \cdot 599 \right) = 359.
\end{align*}

\begin{figure}[H]
    \centering
    \begin{tikzpicture}
        \def\ex{5.0}
        \begin{scope}
            \resetpainel \presetmundo
            \def\pw{3.3}\def\ph{2.3}\def\lfont{\scriptsize}
            \def\lxa{$-2$}\def\lxb{$6$}\def\lya{$-1$}\def\lyb{$4$}\def\lpt{$P_w = (4,\ 2)$}
            \painel{A}
        \end{scope}
        \begin{scope}[xshift=\ex cm]
            \resetpainel \presetndcii
            \def\pw{3.3}\def\ph{2.3}\def\ptit{NDC $[-1,1]^2$}
            \def\lpt{$P_n = (0{,}5;\ 0{,}2)$}
            \painel{B}
        \end{scope}
        \begin{scope}[xshift=2*\ex cm]
            \resetpainel \presetdc
            \def\pw{3.3}\def\ph{2.3}\def\ptit{DC}\def\lfont{\scriptsize}
            \def\lxb{$799$}\def\lyb{$599$}\def\lpt{$P_d = (599,\ 359)$}
            \painel{C}
        \end{scope}
        \draw[-{Stealth[length=2.5mm]}, thick, red!75!black] (A-P) -- (B-P);
        \draw[-{Stealth[length=2.5mm]}, thick, red!75!black] (B-P) -- (C-P);
    \end{tikzpicture}
    \caption{Exemplo numérico do cenário (ii): o mesmo ponto no mundo, no NDC $[-1,1]$ e no dispositivo. A ida percorre os sistemas da esquerda para a direita e a volta, da direita para a esquerda.}
    \label{fig:exemplo-ii}
\end{figure}

\paragraph{Volta (DC $\rightarrow$ NDC $\rightarrow$ Mundo).} Partindo de $P_d = (599, 359)$, pelas Equações~\eqref{eq:dc-ndc-i} e~\eqref{eq:ndc-user-i} (cenário (i)):
\begin{align*}
    x_n^{(i)} &= \frac{599}{799} \approx 0{,}7497, &
    y_n^{(i)} &= \frac{359}{599} \approx 0{,}5993, \\
    x_w &= -2 + \frac{599}{799} \cdot 8 \approx 3{,}9975, &
    y_w &= -1 + \frac{359}{599} \cdot 5 \approx 1{,}9967.
\end{align*}
Pelas Equações~\eqref{eq:dc-ndc-ii} e~\eqref{eq:ndc-user-ii} (cenário (ii)):
\begin{align*}
    x_n^{(ii)} &= 2 \cdot \frac{599}{799} - 1 \approx 0{,}4994, &
    y_n^{(ii)} &= 2 \cdot \frac{359}{599} - 1 \approx 0{,}1987, \\
    x_w &= -2 + \frac{0{,}4994 + 1}{2} \cdot 8 \approx 3{,}9975, &
    y_w &= -1 + \frac{0{,}1987 + 1}{2} \cdot 5 \approx 1{,}9967.
\end{align*}
Em ambos os cenários o ponto recuperado, $P_w \approx (3{,}9975;\ 1{,}9967)$, difere do original $(4,2)$ por menos de meio pixel (cerca de $0{,}005$ em $x$ e $0{,}004$ em $y$), como esperado pelo arredondamento da ida.

\section{Metodologia e desenvolvimento da solução}
\label{sec:metodologia}

\subsection{Arquitetura do programa}
A solução proposta foi desenvolvida em Java, com uso das bibliotecas gráficas nativas da linguagem (Swing e AWT). A arquitetura é composta pelas classes \texttt{Main}, \texttt{JanelaPrincipal}, \texttt{Ponto}, \texttt{JanelaMundo}, \texttt{Transformacao} e \texttt{Display}, e pela enumeração \texttt{TipoNDC}.

A classe \texttt{Main} é responsável pela inicialização do programa. A \texttt{JanelaPrincipal} concentra a interface gráfica e coordena a interação com o usuário, a classe \texttt{Ponto} representa coordenadas, \texttt{JanelaMundo} armazena os limites definidos para o espaço de mundo, a classe \texttt{Transformacao} concentra as transformações utilizadas durante o programa, já a classe \texttt{Display} é responsável pela representação gráfica e pela ativação do pixel correspondente à posição calculada. A enumeração \texttt{TipoNDC} identifica os dois sistemas normalizados disponíveis, \texttt{ZERO\_A\_UM} (NDC \mbox{$[0,1] \times [0,1]$}) e \texttt{MENOS\_UM\_A\_UM} (NDC \mbox{$[-1,1] \times [-1,1]$}).

\subsection{Entradas do sistema}
A aplicação aceita duas formas de entrada, que percorrem o pipeline de transformações em sentidos diferentes.

\paragraph{Clique do mouse.} A primeira entrada é a posição de um clique do mouse sobre o display. O display possui resolução de $800 \times 600$ pixels, sendo que a origem das coordenadas começa no canto superior esquerdo.

Quando o usuário clica sobre o display, a aplicação obtém as coordenadas $x$ e $y$ fornecidas pelo evento de clique do mouse e cria um objeto do tipo \texttt{Ponto} com tais valores. A coordenada obtida, $P_{inp} = (x_{inp}, y_{inp})$, é utilizada como entrada para as transformações. O clique é capturado por um \textit{listener} associado ao display, de modo que cada clique substitui o pixel ativo anterior pela nova posição.

\paragraph{Consulta de ponto do mundo.} A segunda entrada é um ponto $P_w = (x_w, y_w)$ digitado pelo usuário no painel \textit{Consultar Ponto do Mundo}. Ao pressionar o botão \textit{Mostrar no display}, o ponto é convertido para coordenadas de dispositivo e o pixel correspondente é ativado. Essa entrada permite verificar diretamente o mapeamento Mundo $\rightarrow$ DC para valores escolhidos, como os do exemplo numérico da Seção~\ref{sec:fundamentacao-teorica}.

Em ambos os casos, a interface apresenta os valores correspondentes à coordenada de dispositivo, ao NDC \mbox{[0,1]}, ao NDC \mbox{[-1,1]} e à posição equivalente na janela de mundo.

\subsection{Seleção do sistema NDC}
\label{subsec:selecao-ndc}
Um seletor (\textit{combo box}) no painel \textit{Sistema NDC} permite escolher qual dos dois cenários de normalização é utilizado no pipeline: \texttt{ZERO\_A\_UM}, correspondente ao cenário (i), ou \texttt{MENOS\_UM\_A\_UM}, correspondente ao cenário (ii). O valor inicial é \texttt{ZERO\_A\_UM}.

A troca do sistema altera apenas quais funções de transformação serão utilizadas nas próximas interações; ela não dispara, por si só, uma nova ativação de pixel. Dessa forma, o pixel não é deslocado para um ponto que o usuário não solicitou, e o próximo clique ou consulta já utiliza o sistema escolhido. Independentemente da seleção, a interface sempre exibe o ponto nos dois sistemas NDC, o que permite compará-los.

\subsection{Transformações implementadas}
As transformações foram implementadas na classe \texttt{Transformacao}, aplicando diretamente as expressões já deduzidas na Seção~\ref{sec:fundamentacao-teorica} (Fundamentação teórica), que coincidem com as fórmulas da \aularef\ (inclusive por não inverterem o eixo $y$ entre o NDC e o dispositivo), para $W$ e $H$ correspondendo à largura e à altura do display. Quando o sistema selecionado é o NDC \mbox{$[0,1] \times [0,1]$}, o caminho utilizado para a ativação do pixel emprega as seguintes expressões:

\begin{itemize}
    \item \textbf{Entrada $\rightarrow$ NDC} (\texttt{inpToNdc01}), mesma forma da Equação~\eqref{eq:dc-ndc-i}, com $x_{inp}, y_{inp}$ no lugar de $x_d, y_d$:
    \[
        x_n^{(i)} = \frac{x_{inp}}{W-1}, \qquad y_n^{(i)} = \frac{y_{inp}}{H-1}.
    \]
    \item \textbf{NDC $\rightarrow$ Mundo} (\texttt{ndc01ToUser}), Equação~\eqref{eq:ndc-user-i}:
    \[
        x_w = x_{wmin} + x_n^{(i)} \, (x_{wmax} - x_{wmin}), \qquad
        y_w = y_{wmin} + y_n^{(i)} \, (y_{wmax} - y_{wmin}).
    \]
    \item \textbf{Mundo $\rightarrow$ NDC} (\texttt{userToNdc01}), Equação~\eqref{eq:user-ndc-i}:
    \[
        x_n^{(i)} = \frac{x_w - x_{wmin}}{x_{wmax} - x_{wmin}}, \qquad
        y_n^{(i)} = \frac{y_w - y_{wmin}}{y_{wmax} - y_{wmin}}.
    \]
    \item \textbf{NDC $\rightarrow$ DC} (\texttt{ndc01ToDc}), Equação~\eqref{eq:ndc-dc-i} sem o arredondamento final, aplicado apenas na etapa de ativação do pixel (Seção~\ref{subsec:ativacao-pixel}):
    \[
        x_d = x_n^{(i)} \, (W-1), \qquad y_d = y_n^{(i)} \, (H-1).
    \]
\end{itemize}

Quando o sistema selecionado é o NDC \mbox{$[-1,1] \times [-1,1]$}, as mesmas quatro etapas são realizadas pelos métodos \texttt{inpToNdc11}, \texttt{ndc11ToUser}, \texttt{userToNdc11} e \texttt{ndc11ToDc}, conforme as Equações~\eqref{eq:dc-ndc-ii}, \eqref{eq:ndc-user-ii}, \eqref{eq:user-ndc-ii} e \eqref{eq:ndc-dc-ii}. Como demonstrado na Subseção~\ref{subsec:sintese}, os dois cenários levam ao mesmo pixel final; a escolha altera apenas os valores intermediários exibidos e o conjunto de funções exercitado.

Vale observar que a implementação não inverte o eixo $y$ entre o NDC e o dispositivo, como nas fórmulas da aula, e que o raster do Java tem origem no canto superior esquerdo, com $y$ crescendo para baixo. Assim, diferentemente da convenção adotada na Seção~\ref{sec:fundamentacao-teorica} (origem no canto inferior esquerdo, $y$ para cima), no display um ponto com $y_w$ maior é ativado em uma linha mais baixa, e a imagem aparece espelhada verticalmente.

\subsection{Ativação do pixel}
\label{subsec:ativacao-pixel}
Após a realização das transformações, as coordenadas finais do dispositivo podem possuir valores fracionários. Como a posição de um pixel é representada por coordenadas inteiras, os valores obtidos são arredondados utilizando \texttt{Math.round()}. As coordenadas resultantes são enviadas ao método \texttt{drawPixel()} pertencente à classe \texttt{Display}.

A classe utiliza um \texttt{BufferedImage} com resolução de $800 \times 600$ pixels e fundo preto. O método \texttt{drawPixel()} verifica se a posição está dentro dos limites do display e, caso esteja, altera diretamente a cor daquele pixel por meio do método \texttt{setRGB()}. O pixel ativo é representado pela cor verde (\texttt{\#00FF00}). Caso a posição esteja fora dos limites, o display apenas é limpo e nenhum pixel é ativado.

Antes de ativar o novo pixel, o display é limpo. Dessa forma, somente a posição correspondente à coordenada atual permanece destacada.

\subsection{Interface para definição da janela de mundo}
A aplicação possui uma pequena interface que permite ao usuário definir os limites da janela de mundo. São disponibilizados campos para:
\begin{itemize}
    \item X mínimo ($x_{wmin}$);
    \item X máximo ($x_{wmax}$);
    \item Y mínimo ($y_{wmin}$);
    \item Y máximo ($y_{wmax}$).
\end{itemize}
Inicialmente, os campos são preenchidos com os valores $x_{wmin} = -100$, $x_{wmax} = 100$, $y_{wmin} = -100$ e $y_{wmax} = 100$.

O usuário pode alterar esses valores e pressionar o botão \textit{Aplicar}. A aplicação realiza a conversão dos valores informados para o tipo \texttt{double} e cria uma nova instância de \texttt{JanelaMundo} com os limites fornecidos. A nova janela passa a ser utilizada a partir da próxima interação (clique ou consulta).

A classe \texttt{JanelaMundo} verifica se os valores fornecidos formam uma janela válida, exigindo que o valor mínimo seja menor que o valor máximo em ambos os eixos. Caso sejam informados valores não numéricos ou limites inválidos, a aplicação apresenta uma mensagem de erro ao usuário e mantém a janela anterior.

\subsection{Fluxo geral da solução}
De forma resumida, o funcionamento da solução pode ser representado por dois fluxos. O primeiro corresponde ao clique do mouse:
\[
    \text{Mouse} \Rightarrow \text{INP} \Rightarrow \text{NDC} \Rightarrow \text{USER} \Rightarrow \text{NDC} \Rightarrow \text{DC} \Rightarrow \text{Pixel}
\]
O mouse fornece a coordenada inicial utilizada pelo nosso programa, as funções da classe \texttt{Transformacao} realizam as conversões, a janela de mundo fornece os limites para a transformação, e o \texttt{Display} utiliza a coordenada final para ativar diretamente o pixel correspondente. O percurso completo, de ida ao mundo e volta ao dispositivo, garante que o pixel ativado seja o resultado das transformações no sistema NDC selecionado, e não simplesmente a posição bruta do clique.

O segundo fluxo corresponde à consulta de um ponto do mundo:
\[
    \text{USER} \Rightarrow \text{NDC} \Rightarrow \text{DC} \Rightarrow \text{Pixel}
\]
Em ambos os fluxos, a etapa NDC utiliza o sistema escolhido no seletor (Subseção~\ref{subsec:selecao-ndc}).

A aplicação apresenta as coordenadas intermediárias na interface, permitindo acompanhar visualmente o resultado das transformações realizadas.

\section{Implementação}

\subsection{Organização do código}
A implementação foi estruturada em diferentes classes, de modo a separar as responsabilidades da aplicação. Essa organização permite que a representação dos dados, os cálculos matemáticos, a interface gráfica e a manipulação do display sejam tratados de forma independente.

A classe \texttt{Ponto} foi utilizada para representar coordenadas bidimensionais. Ela possui os atributos \texttt{x} e \texttt{y}, além de métodos de acesso e alteração desses valores. Dessa forma, as coordenadas utilizadas durante as transformações podem ser encapsuladas em objetos.

A classe \texttt{JanelaMundo} concentra os valores referentes aos limites do espaço de mundo. Seu construtor e seus métodos de alteração realizam verificações para garantir que os valores mínimos sejam menores que os valores máximos.

A classe \texttt{Transformacao} concentra os métodos responsáveis pelos cálculos matemáticos. Os métodos recebem os pontos e os parâmetros necessários para cada transformação e retornam novos objetos \texttt{Ponto} contendo os valores calculados.

A classe \texttt{Display} é responsável pela representação do display por meio de um \texttt{BufferedImage}. Já a classe \texttt{JanelaPrincipal} realiza a integração entre os componentes, controlando a interface, a entrada do mouse e a execução das transformações.

Por fim, a classe \texttt{Main} contém o ponto de inicialização da aplicação e cria a janela principal.

\subsection{Implementação das funções de transformação}
As funções de transformação foram implementadas de forma independente na classe \texttt{Transformacao}. Essa abordagem permite que cada operação matemática seja executada por um método específico.

Para a conversão da entrada para NDC \mbox{[0,1]}, foi implementado o método \texttt{inpToNdc01()}. Para NDC \mbox{[-1,1]}, foi utilizado o método \texttt{inpToNdc11()}.

As transformações entre NDC e a janela de mundo foram separadas nos métodos \texttt{ndc01ToUser()}, \texttt{userToNdc01()}, \texttt{ndc11ToUser()} e \texttt{userToNdc11()}. Também foram implementados os métodos \texttt{ndc01ToDc()} e \texttt{ndc11ToDc()}, responsáveis pela conversão das coordenadas normalizadas para as coordenadas do dispositivo.

Além dos cálculos, a classe possui o método privado \texttt{validarDimensoes()}, utilizado para verificar se as dimensões fornecidas para o display são válidas antes da realização das operações que dependem da largura e da altura.

\subsection{Implementação da entrada do mouse}
O rastreamento do mouse foi implementado na classe \texttt{JanelaPrincipal}, no método \texttt{configurarMouse()}, por meio de um \texttt{MouseInputAdapter} que sobrescreve três métodos: \texttt{mousePressed()}, chamado quando um botão é pressionado; \texttt{mouseMoved()}, chamado a cada movimento do mouse sobre o display; e \texttt{mouseDragged()}, chamado a cada movimento com o botão pressionado. O mesmo adaptador é registrado no display duas vezes, com \texttt{addMouseListener()} e com \texttt{addMouseMotionListener()}. Esse duplo registro é necessário porque o Swing entrega \texttt{mousePressed()} apenas aos \texttt{MouseListener}, enquanto \texttt{mouseMoved()} e \texttt{mouseDragged()} são entregues apenas aos \texttt{MouseMotionListener}.

Os três métodos encaminham o evento ao método \texttt{rastrearMouse()}, que obtém as coordenadas em relação ao display e verifica, por meio de \texttt{dentroDoDisplay()}, se elas estão dentro dos limites. Em caso positivo, os valores são encaminhados ao método \texttt{ativarPixelPorClique()}, que cria um objeto do tipo \texttt{Ponto} e o repassa aos métodos da classe \texttt{Transformacao}. Apesar do nome, mantido da versão inicial em que apenas o clique era tratado, esse método é executado para qualquer um dos três eventos.

Em \texttt{ativarPixelPorClique()}, o ponto de entrada é convertido para os dois sistemas NDC (\texttt{inpToNdc01()} e \texttt{inpToNdc11()}). Em seguida, o ponto do mundo é obtido a partir do NDC selecionado (\texttt{ndc01ToUser()} ou \texttt{ndc11ToUser()}) e convertido de volta para coordenadas de dispositivo pelo método auxiliar \texttt{mundoParaDispositivo()}. A atualização do pixel ativo e das informações da interface ocorre ao final desse método, de modo que cada nova posição do mouse produz uma nova ativação.

\subsection{Implementação da seleção do sistema NDC}
O sistema NDC em uso é armazenado no atributo \texttt{ndcAtual} da \texttt{JanelaPrincipal}, do tipo \texttt{TipoNDC}. O seletor é um \texttt{JComboBox} preenchido com \texttt{TipoNDC.values()}; seu \textit{listener} apenas atualiza \texttt{ndcAtual} com o item escolhido.

O método \texttt{mundoParaDispositivo()} centraliza o uso desse atributo: quando \texttt{ndcAtual} vale \texttt{ZERO\_A\_UM}, utiliza \texttt{userToNdc01()} seguido de \texttt{ndc01ToDc()}; caso contrário, utiliza \texttt{userToNdc11()} seguido de \texttt{ndc11ToDc()}. Esse método é compartilhado pelas duas formas de entrada, o que garante que o mouse e a consulta de ponto empreguem sempre o mesmo pipeline.

\subsection{Implementação da consulta de ponto do mundo}
O painel \textit{Consultar Ponto do Mundo} possui dois campos de texto, para $x_w$ e $y_w$, e o botão \textit{Mostrar no display}. Ao ser pressionado, o botão executa o método \texttt{consultarPontoMundo()}, que converte os valores digitados para \texttt{double}, cria o \texttt{Ponto} correspondente e obtém suas coordenadas de dispositivo por meio de \texttt{mundoParaDispositivo()}. O resultado é arredondado com \texttt{Math.round()} e, se estiver dentro do display, enviado a \texttt{drawPixel()}; caso contrário, o pixel não é ativado e um alerta é exibido (Subseção~\ref{subsec:impl-fora-limites}). Valores não numéricos são tratados por uma \texttt{NumberFormatException}, que resulta em uma mensagem de erro.

Os valores exibidos na parte inferior da janela são atualizados pelo método \texttt{atualizarLabels()}. Na consulta de ponto, esse método recalcula o ponto nos dois sistemas NDC (\texttt{userToNdc01()} e \texttt{userToNdc11()}) apenas para fins de exibição. Ao ser chamado, o método também apaga a mensagem de alerta que estiver sendo exibida, de modo que o alerta desaparece assim que uma coordenada válida volta a ser informada.

\subsection{Implementação do tratamento de coordenadas fora dos limites}
\label{subsec:impl-fora-limites}
A verificação dos limites é feita pelo método \texttt{dentroDoDisplay()}, que retorna verdadeiro apenas quando $0 \le x < 800$ e $0 \le y < 600$. Ela é aplicada em dois pontos: em \texttt{rastrearMouse()}, sobre a posição do mouse, e em \texttt{consultarPontoMundo()}, sobre o pixel obtido para o ponto do mundo, já arredondado.

Quando a coordenada está fora dos limites, é chamado o método \texttt{mostrarForaDosLimites()}, que limpa o display, substitui os textos de dispositivo, NDC \mbox{[0,1]}, NDC \mbox{[-1,1]} e mundo por ``fora dos limites'' e escreve a mensagem recebida em um \texttt{JLabel} de alerta, exibido em vermelho abaixo das coordenadas. Para o mouse, a mensagem informa a posição e a resolução do display, por exemplo ``Coordenada ($-15$, 320) fora dos limites do display 800x600.''. Para a consulta de ponto, a mensagem informa o ponto digitado e os limites da janela de mundo e também é apresentada em uma caixa de diálogo de aviso (\texttt{JOptionPane.WARNING\_MESSAGE}).

\subsection{Implementação do display}
A representação do display foi implementada na classe \texttt{Display} utilizando a classe \texttt{BufferedImage}. A imagem possui exatamente $800 \times 600$ posições, correspondentes à resolução que foi definida para a aplicação.

O método \texttt{drawPixel()} recebe as coordenadas inteiras do pixel que deverá ser ativado. Antes de realizar a alteração, o método verifica se as coordenadas estão dentro dos limites da imagem. Quando a posição é válida, o método \texttt{setRGB()} altera diretamente o valor de cor daquela posição. A cor utilizada para o pixel ativo é obtida a partir de \texttt{Color.GREEN}.

A limpeza do display é realizada pelo método \texttt{limpar()}, que preenche a imagem com a cor de fundo antes da ativação de uma nova posição. A imagem armazenada no \texttt{BufferedImage} é posteriormente apresentada pelo método \texttt{paintComponent()}, responsável por desenhá-la no componente gráfico.

\subsection{Implementação da interface e configuração da janela de mundo}
A interface gráfica foi construída utilizando componentes do Swing. A classe \texttt{JanelaPrincipal} utiliza um \texttt{BorderLayout}: o display ocupa a região central, as informações das coordenadas e a linha de alerta ficam na região inferior e um painel lateral, à direita, agrupa os painéis \textit{Janela do Mundo}, \textit{Sistema NDC} e \textit{Consultar Ponto do Mundo}. A janela é dimensionada com \texttt{pack()}, centralizada na tela e não pode ser redimensionada, o que mantém o display com exatamente $800 \times 600$ pixels.

Os campos de texto são utilizados para receber os valores de $x_{wmin}$, $x_{wmax}$, $y_{wmin}$ e $y_{wmax}$. Quando o usuário solicita a aplicação dos novos valores, o método responsável pela atualização realiza a conversão dos dados para \texttt{double} e cria uma nova instância de \texttt{JanelaMundo}.

A implementação também trata possíveis entradas inválidas. Valores que não podem ser convertidos para números ou que não respeitam os limites necessários são tratados por exceções e apresentados ao usuário por meio de mensagens na interface.

\subsection{Integração da implementação}
A integração entre as classes ocorre principalmente na \texttt{JanelaPrincipal}. Quando a posição do mouse muda ou um ponto do mundo é consultado, a classe cria o ponto inicial e utiliza os métodos da \texttt{Transformacao}, de acordo com o sistema NDC selecionado, para obter os valores correspondentes aos diferentes sistemas de coordenadas. Depois das transformações necessárias, a coordenada final é arredondada com \texttt{Math.round()}, verificada em relação aos limites do display e, se válida, enviada para o método \texttt{drawPixel()} da classe \texttt{Display}.

A separação entre essas responsabilidades permite que os cálculos de transformação permaneçam independentes da implementação gráfica. Ao adotar esse modelo, a classe \texttt{Transformacao} não precisa conhecer os componentes da interface ou a forma como o pixel é desenhado, enquanto \texttt{Display} não precisa realizar os cálculos matemáticos das coordenadas.

Essa estrutura mantém a implementação organizada e permite relacionar diretamente cada parte do código com sua respectiva função dentro da aplicação.

\section{Resultados}
\label{sec:resultados}

Esta seção apresenta os resultados obtidos com a execução do programa. Os valores numéricos das tabelas foram obtidos por meio de uma simulação que utilizou as próprias classes \texttt{Ponto}, \texttt{JanelaMundo}, \texttt{TipoNDC} e \texttt{Transformacao} do projeto, sem modificações, reproduzindo a mesma sequência de chamadas realizada pela \texttt{JanelaPrincipal} nos métodos \texttt{ativarPixelPorClique()}, \texttt{consultarPontoMundo()} e \texttt{mundoParaDispositivo()}, com um display de $W = 800$ por $H = 600$ pixels. As capturas de tela reproduzem, na interface, os mesmos casos da simulação. Como o pixel ativo ocupa uma única posição do display, ele foi destacado nas capturas por um círculo vermelho, acrescentado à imagem apenas para facilitar sua localização.

Nas tabelas, os pares de coordenadas reais são escritos com ponto e vírgula entre as componentes, e as coordenadas de dispositivo são apresentadas antes do arredondamento (\emph{DC real}) e depois dele (\emph{pixel}).

% ------------------------------------------------------------
\subsection{Funcionamento do programa}
Ao ser iniciado, o programa exibe o display de $800 \times 600$ pixels com fundo preto, o painel lateral com os campos da janela de mundo (inicialmente $[-100, 100] \times [-100, 100]$), o seletor do sistema NDC e o painel de consulta de ponto, além da área inferior em que são mostradas as coordenadas de dispositivo, NDC \mbox{[0,1]}, NDC \mbox{[-1,1]} e mundo.

As Figuras~\ref{fig:resultado-ndc01} e~\ref{fig:resultado-ndc11} mostram a consulta do ponto $P_w = (0, 0)$, centro da janela de mundo padrão, com cada um dos sistemas NDC selecionados. Em ambos os casos, o pixel ativado é $(400, 300)$ e os valores exibidos são idênticos, pois a interface sempre apresenta o ponto nos dois sistemas; muda apenas o conjunto de funções utilizado internamente para chegar ao pixel.

\begin{figure}[H]
    \centering
    \includegraphics[width=0.75\textwidth]{ndc01.png}
    \caption{Consulta do ponto $P_w = (0, 0)$ com o sistema NDC $[0,1]$ (\texttt{ZERO\_A\_UM}) selecionado.}
    \label{fig:resultado-ndc01}
\end{figure}

\begin{figure}[H]
    \centering
    \includegraphics[width=0.75\textwidth]{ndc11.png}
    \caption{Consulta do mesmo ponto $P_w = (0, 0)$ com o sistema NDC $[-1,1]$ (\texttt{MENOS\_UM\_A\_UM}) selecionado. O pixel ativado é o mesmo da Figura~\ref{fig:resultado-ndc01}.}
    \label{fig:resultado-ndc11}
\end{figure}

O valor $(400, 300)$ decorre do arredondamento: o centro exato do display é $(399{,}5;\ 299{,}5)$, pois os índices dos pixels vão de $0$ a $W-1 = 799$ e de $0$ a $H-1 = 599$, e a função \texttt{Math.round()} arredonda a metade para cima.

% ------------------------------------------------------------
\subsection{Consulta de pontos do mundo}
\label{subsec:res-consulta}
Com a janela de mundo padrão, foram consultados os quatro cantos da janela, seu centro e alguns pontos internos. A Tabela~\ref{tab:res-consulta} apresenta os resultados. Para todos os pontos, os dois sistemas NDC produziram exatamente o mesmo pixel, de modo que apenas uma coluna de pixel é apresentada.

\begin{table}[H]
    \centering
    \renewcommand{\arraystretch}{1.35}
    \small
    \begin{tabular}{c|c|c|c|c}
        \textbf{Mundo} & \textbf{NDC [0,1]} & \textbf{NDC [-1,1]} & \textbf{DC real} & \textbf{Pixel} \\ \hline
        $(0;\ 0)$           & $(0{,}5;\ 0{,}5)$        & $(0;\ 0)$                 & $(399{,}50;\ 299{,}50)$ & $(400, 300)$ \\
        $(-100;\ -100)$     & $(0;\ 0)$                & $(-1;\ -1)$               & $(0;\ 0)$               & $(0, 0)$ \\
        $(100;\ 100)$       & $(1;\ 1)$                & $(1;\ 1)$                 & $(799;\ 599)$           & $(799, 599)$ \\
        $(-100;\ 100)$      & $(0;\ 1)$                & $(-1;\ 1)$                & $(0;\ 599)$             & $(0, 599)$ \\
        $(100;\ -100)$      & $(1;\ 0)$                & $(1;\ -1)$                & $(799;\ 0)$             & $(799, 0)$ \\
        $(50;\ 50)$         & $(0{,}75;\ 0{,}75)$      & $(0{,}5;\ 0{,}5)$         & $(599{,}25;\ 449{,}25)$ & $(599, 449)$ \\
        $(-50;\ 25)$        & $(0{,}25;\ 0{,}625)$     & $(-0{,}5;\ 0{,}25)$       & $(199{,}75;\ 374{,}38)$ & $(200, 374)$ \\
        $(33{,}3;\ -66{,}6)$ & $(0{,}6665;\ 0{,}167)$  & $(0{,}333;\ -0{,}666)$    & $(532{,}53;\ 100{,}03)$ & $(533, 100)$
    \end{tabular}
    \caption{Consulta de pontos do mundo com a janela $[-100, 100] \times [-100, 100]$.}
    \label{tab:res-consulta}
\end{table}

Os cantos da janela de mundo foram levados exatamente aos cantos do display, e os valores normalizados permaneceram nos intervalos esperados, $[0,1]$ e $[-1,1]$, para todos os pontos. Os valores do NDC $[-1,1]$ obedecem à relação $x_n^{(ii)} = 2\,x_n^{(i)} - 1$ em todas as linhas, como previsto na Subseção~\ref{subsec:sintese}.

\begin{figure}[H]
    \centering
    \includegraphics[width=0.75\textwidth]{consulta50-50.png}
    \caption{Consulta do ponto $P_w = (50, 50)$ com a janela de mundo padrão.}
    \label{fig:res-consulta-50-50}
\end{figure}

% ------------------------------------------------------------
\subsection{Ativação de pixels por clique}
\label{subsec:res-clique}
A Tabela~\ref{tab:res-clique} apresenta o resultado de cliques em diferentes posições do display, com a janela de mundo padrão. Para cada clique, o ponto foi levado ao mundo pelo NDC selecionado e trazido de volta ao dispositivo. As colunas de mundo e pixel final são as mesmas nos dois sistemas NDC.

\begin{table}[H]
    \centering
    \renewcommand{\arraystretch}{1.35}
    \small
    \begin{tabular}{c|c|c|c|c}
        \textbf{Clique} & \textbf{NDC [0,1]} & \textbf{NDC [-1,1]} & \textbf{Mundo} & \textbf{Pixel final} \\ \hline
        $(0, 0)$     & $(0;\ 0)$                 & $(-1;\ -1)$                 & $(-100;\ -100)$             & $(0, 0)$ \\
        $(799, 0)$   & $(1;\ 0)$                 & $(1;\ -1)$                  & $(100;\ -100)$              & $(799, 0)$ \\
        $(0, 599)$   & $(0;\ 1)$                 & $(-1;\ 1)$                  & $(-100;\ 100)$              & $(0, 599)$ \\
        $(799, 599)$ & $(1;\ 1)$                 & $(1;\ 1)$                   & $(100;\ 100)$               & $(799, 599)$ \\
        $(400, 300)$ & $(0{,}5006;\ 0{,}5008)$   & $(0{,}0013;\ 0{,}0017)$     & $(0{,}1252;\ 0{,}1669)$     & $(400, 300)$ \\
        $(200, 150)$ & $(0{,}2503;\ 0{,}2504)$   & $(-0{,}4994;\ -0{,}4992)$   & $(-49{,}9374;\ -49{,}9165)$ & $(200, 150)$ \\
        $(600, 450)$ & $(0{,}7509;\ 0{,}7513)$   & $(0{,}5019;\ 0{,}5025)$     & $(50{,}1877;\ 50{,}2504)$   & $(600, 450)$ \\
        $(123, 456)$ & $(0{,}1539;\ 0{,}7613)$   & $(-0{,}6921;\ 0{,}5225)$    & $(-69{,}2115;\ 52{,}2538)$  & $(123, 456)$
    \end{tabular}
    \caption{Cliques no display com a janela $[-100, 100] \times [-100, 100]$.}
    \label{tab:res-clique}
\end{table}

Em todos os casos, o pixel final coincidiu com o pixel clicado, o que mostra que o percurso de ida e volta pelo mundo não desloca o ponto. Os cliques nos cantos do display correspondem exatamente aos cantos da janela de mundo.

O clique em $(400, 300)$ não corresponde exatamente à origem do mundo, mas a $(0{,}1252;\ 0{,}1669)$. Isso ocorre porque o centro do display, $(399{,}5;\ 299{,}5)$, não é um pixel inteiro, e cada pixel representa uma faixa de valores do mundo. Com a janela padrão, um pixel equivale a $200/799 \approx 0{,}2503$ unidades do mundo na horizontal e a $200/599 \approx 0{,}3339$ unidades na vertical, de modo que o deslocamento observado é de meio pixel em cada eixo.

\begin{figure}[H]
    \centering
    \includegraphics[width=0.75\textwidth]{clique-ndc01.png}
    \caption{Ativação de pixel por clique na posição $(200, 150)$ com o sistema NDC $[0,1]$ selecionado.}
    \label{fig:res-clique-ndc01}
\end{figure}

\begin{figure}[H]
    \centering
    \includegraphics[width=0.75\textwidth]{clique-ndc11.png}
    \caption{Ativação de pixel por clique na mesma posição, $(200, 150)$, com o sistema NDC $[-1,1]$ selecionado. O pixel e os valores exibidos são os mesmos da Figura~\ref{fig:res-clique-ndc01}.}
    \label{fig:res-clique-ndc11}
\end{figure}

% ------------------------------------------------------------
\subsection{Alteração da janela de mundo}
\label{subsec:res-janela}
Para verificar a alteração da janela de mundo, foi aplicada a janela $[-2, 6] \times [-1, 4]$, a mesma do exemplo numérico da Seção~\ref{sec:fundamentacao-teorica}, e, em seguida, a janela $[0, 1000] \times [0, 10]$, cujos eixos têm escalas muito diferentes. A Tabela~\ref{tab:res-janela} apresenta os resultados.

\begin{table}[H]
    \centering
    \renewcommand{\arraystretch}{1.35}
    \small
    \begin{tabular}{c|c|c|c|c}
        \textbf{Entrada} & \textbf{NDC [0,1]} & \textbf{NDC [-1,1]} & \textbf{Mundo} & \textbf{Pixel} \\ \hline
        \multicolumn{5}{c}{Janela $[-2, 6] \times [-1, 4]$} \\ \hline
        consulta $(4;\ 2)$    & $(0{,}75;\ 0{,}6)$       & $(0{,}5;\ 0{,}2)$         & $(4;\ 2)$                 & $(599, 359)$ \\
        consulta $(0;\ 0)$    & $(0{,}25;\ 0{,}2)$       & $(-0{,}5;\ -0{,}6)$       & $(0;\ 0)$                 & $(200, 120)$ \\
        consulta $(-2;\ -1)$  & $(0;\ 0)$                & $(-1;\ -1)$               & $(-2;\ -1)$               & $(0, 0)$ \\
        consulta $(6;\ 4)$    & $(1;\ 1)$                & $(1;\ 1)$                 & $(6;\ 4)$                 & $(799, 599)$ \\
        clique $(599, 359)$   & $(0{,}7497;\ 0{,}5993)$  & $(0{,}4994;\ 0{,}1987)$   & $(3{,}9975;\ 1{,}9967)$   & $(599, 359)$ \\
        clique $(400, 300)$   & $(0{,}5006;\ 0{,}5008)$  & $(0{,}0013;\ 0{,}0017)$   & $(2{,}0050;\ 1{,}5042)$   & $(400, 300)$ \\ \hline
        \multicolumn{5}{c}{Janela $[0, 1000] \times [0, 10]$} \\ \hline
        consulta $(500;\ 5)$    & $(0{,}5;\ 0{,}5)$      & $(0;\ 0)$                 & $(500;\ 5)$               & $(400, 300)$ \\
        consulta $(250;\ 7{,}5)$ & $(0{,}25;\ 0{,}75)$   & $(-0{,}5;\ 0{,}5)$        & $(250;\ 7{,}5)$           & $(200, 449)$ \\
        clique $(400, 300)$     & $(0{,}5006;\ 0{,}5008)$ & $(0{,}0013;\ 0{,}0017)$  & $(500{,}6258;\ 5{,}0083)$ & $(400, 300)$
    \end{tabular}
    \caption{Resultados após a alteração da janela de mundo.}
    \label{tab:res-janela}
\end{table}

Com a janela $[-2, 6] \times [-1, 4]$, a consulta de $P_w = (4, 2)$ ativou o pixel $(599, 359)$, e o clique nesse pixel recuperou o ponto $(3{,}9975;\ 1{,}9967)$. Esses valores coincidem com os calculados manualmente no exemplo numérico da Seção~\ref{sec:fundamentacao-teorica} (Figuras~\ref{fig:exemplo-i} e~\ref{fig:exemplo-ii}), o que valida a implementação em relação às deduções.

Os valores de NDC do clique em $(400, 300)$ são os mesmos para todas as janelas, pois a normalização da entrada depende apenas do tamanho do display. A alteração da janela modifica somente a etapa NDC $\rightarrow$ Mundo. Com a janela $[0, 1000] \times [0, 10]$, cada pixel representa cerca de $1{,}2516$ unidades do mundo na horizontal e apenas $0{,}0167$ na vertical, o que evidencia que a transformação aplica escalas independentes em cada eixo.

\begin{figure}[H]
    \centering
    \includegraphics[width=0.75\textwidth]{janela-exemplo.png}
    \caption{Consulta do ponto do exemplo numérico após a alteração da janela de mundo.}
    \label{fig:res-janela-exemplo}
\end{figure}

\begin{figure}[H]
    \centering
    \includegraphics[width=0.75\textwidth]{janela-assimetrica.png}
    \caption{Consulta de ponto com uma janela de mundo de escalas diferentes em cada eixo.}
    \label{fig:res-janela-assimetrica}
\end{figure}

% ------------------------------------------------------------
\subsection{Equivalência entre os sistemas NDC}
\label{subsec:res-equivalencia}
Para verificar de forma mais ampla os resultados das subseções anteriores, a simulação percorreu todas as posições do display e um grande número de pontos do mundo, para cada uma das três janelas utilizadas ($[-100, 100]^2$, $[-2, 6] \times [-1, 4]$ e $[0, 1000] \times [0, 10]$). Os resultados foram os mesmos para as três janelas:

\begin{itemize}
    \item \textbf{Ida e volta a partir do dispositivo:} para as $800 \times 600 = 480\,000$ posições de clique possíveis, o pixel final coincidiu com o pixel clicado em todos os casos, tanto com o NDC $[0,1]$ quanto com o NDC $[-1,1]$;
    \item \textbf{Mesmo ponto no mundo:} a maior diferença entre o ponto do mundo obtido pelos dois sistemas NDC, para um mesmo clique, foi da ordem de $10^{-14}$, ou seja, apenas erro de arredondamento de ponto flutuante;
    \item \textbf{Mesmo pixel:} para $1\,000\,000$ de pontos do mundo sorteados dentro da janela, os dois sistemas NDC ativaram exatamente o mesmo pixel;
    \item \textbf{Erro de posição:} ao converter esses pontos para pixel e de volta para o mundo, o maior erro observado foi de $0{,}5$ pixel em cada eixo, o limite previsto para o arredondamento ao inteiro mais próximo.
\end{itemize}

Esses resultados confirmam experimentalmente a conclusão da Subseção~\ref{subsec:sintese}: a escolha entre o NDC $[0,1]$ e o NDC $[-1,1]$ altera apenas os valores intermediários, e não o pixel final nem o ponto recuperado no mundo.

% ------------------------------------------------------------
\subsection{Validação das entradas}
\label{subsec:res-validacao}
A Tabela~\ref{tab:res-validacao} apresenta o comportamento do programa para diferentes valores digitados nos campos da janela de mundo. Nos casos de erro, a mensagem é exibida em uma caixa de diálogo e a janela de mundo anterior é mantida.

\begin{table}[H]
    \centering
    \renewcommand{\arraystretch}{1.35}
    \small
    \begin{tabular}{c|c|c|c|l}
        \textbf{X mín.} & \textbf{X máx.} & \textbf{Y mín.} & \textbf{Y máx.} & \textbf{Resultado} \\ \hline
        \texttt{-100} & \texttt{100} & \texttt{-100} & \texttt{100} & janela aceita \\
        \texttt{10}   & \texttt{10}  & \texttt{0}    & \texttt{5}   & ``Xmin deve ser menor que Xmax.'' \\
        \texttt{10}   & \texttt{-10} & \texttt{0}    & \texttt{5}   & ``Xmin deve ser menor que Xmax.'' \\
        \texttt{0}    & \texttt{5}   & \texttt{3}    & \texttt{3}   & ``Ymin deve ser menor que Ymax.'' \\
        \texttt{abc}  & \texttt{1}   & \texttt{0}    & \texttt{1}   & ``Digite apenas valores numéricos.'' \\
        (vazio)       & \texttt{1}   & \texttt{0}    & \texttt{1}   & ``Digite apenas valores numéricos.'' \\
        \texttt{1,5}  & \texttt{3}   & \texttt{0}    & \texttt{1}   & ``Digite apenas valores numéricos.''
    \end{tabular}
    \caption{Validação dos limites da janela de mundo.}
    \label{tab:res-validacao}
\end{table}

A última linha mostra que os valores decimais devem ser digitados com ponto (por exemplo, \texttt{1.5}), pois a conversão é feita por \texttt{Double.parseDouble()}, que não aceita a vírgula como separador decimal. No painel de consulta de ponto, valores não numéricos produzem a mensagem ``Digite apenas valores numéricos para o ponto do mundo.''.

\begin{figure}[H]
    \centering
    \includegraphics[width=0.75\textwidth]{erro-janela.png}
    \caption{Mensagem de erro para uma janela de mundo inválida.}
    \label{fig:res-erro-janela}
\end{figure}

% ------------------------------------------------------------
\subsection{Orientação do eixo \texorpdfstring{$y$}{y}}
\label{subsec:res-orientacao}
Como discutido na Seção~\ref{sec:metodologia}, a implementação não inverte o eixo $y$ e o raster do Java cresce para baixo. A Tabela~\ref{tab:res-consulta} confirma esse comportamento: o ponto $(-100;\ 100)$, canto superior esquerdo da janela de mundo, foi ativado no pixel $(0, 599)$, que corresponde ao canto \emph{inferior} esquerdo do display, enquanto o ponto $(-100;\ -100)$ foi ativado no pixel $(0, 0)$, no canto superior esquerdo. Da mesma forma, o ponto $(-50;\ 25)$, que está acima da origem no mundo, aparece abaixo do centro do display. A imagem no display é, portanto, um espelhamento vertical da imagem no mundo.

\begin{figure}[H]
    \centering
    \includegraphics[width=0.75\textwidth]{orientacao-y.png}
    \caption{Ponto do canto superior esquerdo da janela de mundo ativado no canto inferior esquerdo do display.}
    \label{fig:res-orientacao}
\end{figure}

% ------------------------------------------------------------
\subsection{Síntese dos resultados}
Os testes mostraram que as transformações implementadas reproduzem as expressões deduzidas na Seção~\ref{sec:fundamentacao-teorica}: os cantos e o centro da janela de mundo foram levados às posições esperadas do display, o exemplo numérico calculado manualmente foi reproduzido pelo programa, e os dois sistemas NDC levaram sempre ao mesmo pixel. O percurso de ida e volta a partir de um clique preservou o pixel em todas as posições do display, e o erro de posição ficou limitado a meio pixel, como previsto pelo arredondamento.

Também foram observados comportamentos que devem ser considerados no uso do programa: a orientação do eixo $y$ no display é invertida em relação à convenção da aula e, como a janela padrão é quadrada e o display tem proporção $4:3$, as escalas horizontal e vertical de um pixel no mundo são diferentes.

\section{Conclusão}

Neste laboratório implementamos em Java as transformações entre as coordenadas do mundo, as coordenadas normalizadas e as coordenadas do dispositivo, seguindo as fórmulas vistas em sala. O programa final acompanha a posição do mouse sobre o display o tempo todo, ao mover, clicar ou arrastar, e também permite digitar um ponto do mundo, escolher qual NDC usar e alterar a janela de mundo, mostrando em cada caso os valores de todos os sistemas envolvidos. Quando a coordenada sai dos limites do display, o programa mostra um alerta em vez de valores que não correspondem a nenhum pixel. Com isso, foi possível acompanhar na prática um caminho que, no quadro, parecia apenas uma sequência de regras de três.

Os testes mostraram que a implementação está de acordo com o que foi deduzido. O exemplo que resolvemos à mão na fundamentação, com a janela $[-2, 6] \times [-1, 4]$ e o ponto $(4, 2)$, chegou no programa ao mesmo pixel $(599, 359)$. Também conferimos que o NDC $[0,1]$ e o NDC $[-1,1]$ levam sempre ao mesmo pixel. No começo, esperávamos alguma diferença entre os dois, mas depois de montar as contas ficou claro que eles são só duas formas de escrever a mesma proporção. O mesmo vale para o arredondamento: ao voltar do pixel para o mundo, o ponto nunca é recuperado exatamente, e o erro fica sempre dentro de meio pixel, como era de se esperar.

Algumas questões só apareceram quando o programa começou a rodar. A principal foi a orientação do eixo $y$: no Java a origem fica no canto superior esquerdo e o $y$ cresce para baixo, enquanto em sala usamos a origem embaixo. Como não fizemos a inversão, a imagem no display fica espelhada na vertical. Outra foi a diferença de escala entre os eixos, já que a janela padrão é quadrada e o display não é. Nenhuma das duas altera as contas, mas elas mostram que, além das fórmulas, é preciso levar em conta como o dispositivo realmente organiza os pixels.

Como próximos passos, pensamos em inverter o eixo $y$ na conversão para o dispositivo, para que o display siga a mesma convenção da aula, e em manter a proporção entre os eixos ao definir a janela de mundo. Essas transformações também devem servir de base para os próximos trabalhos da disciplina, em que será necessário desenhar mais do que um único ponto na tela.

\section{Referências}
\begin{thebibliography}{9}

\bibitem{repo}
CG-2026.2. [S. l.]: GitHub, 2026. Disponível em: \url{https://github.com/julioP-dev/CG-2026.2}. Acesso em: 24 set. 2026.

\bibitem{aula2}
SOUSA, Robson Pequeno de. \textit{Sistemas de coordenadas}: Aula 2. Campina Grande: Universidade Estadual da Paraíba, 2026. Material da disciplina de Computação Gráfica.

\bibitem{quadro2}
SOUSA, Robson Pequeno de. \textit{Sistemas de coordenadas}: anotações em quadro da Aula 2 (mundo, NDC e dispositivo desenhados com origem no canto inferior esquerdo e eixo $y$ para cima; deduções de mundo $\rightarrow$ NDC e NDC $\rightarrow$ DC). Campina Grande: Universidade Estadual da Paraíba, 17 set. 2026. Aula expositiva da disciplina de Computação Gráfica, realizada em quinta-feira, das 11h às 13h; registro fotográfico do quadro.

\bibitem{foley}
FOLEY, James D. et al. \textit{Computer graphics}: principles and practice. 2. ed. Reading: Addison-Wesley Professional, 1995.

\bibitem{hearn}
HEARN, Donald; BAKER, M. Pauline. \textit{Computer graphics with OpenGL}. 3. ed. Upper Saddle River: Pearson Prentice Hall, 2004.

\end{thebibliography}


\end{document}
